\documentclass[11pt,a4paper]{article}
\usepackage[T1]{fontenc}
\usepackage[utf8]{inputenc}
\usepackage[indonesian]{babel}
\usepackage{amsmath,amssymb,booktabs,array,graphicx,float}
\usepackage[a4paper,margin=2.2cm]{geometry}
\usepackage{microtype}
% Disable font expansion because the available T1 fonts are not scalable.
\microtypesetup{expansion=false}
\usepackage[hidelinks]{hyperref}
\usepackage{url}
\graphicspath{{Gambar/}}
\setlength{\parskip}{0.35em}
\setlength{\emergencystretch}{2em}
\addto\captionsindonesian{\renewcommand{\figurename}{Gambar}\renewcommand{\tablename}{Tabel}}
\newcommand{\R}{\mathbb R}
\newcommand{\Spec}{\operatorname{Spec}}
\newcommand{\RMS}{\operatorname{RMS}}
\input{hasil_numerik/nilai.tex}
\title{\LARGE Analisis Waktu--Frekuensi Audio Pendamping\\
\textit{Bad Piggies Theme Remake}\\[0.4em]
\large ETS Analisis Fourier dan Wavelet}
\author{Teosofi Hidayah Agung\\6002261009}
\date{22 Oktober 2026}

\begin{document}
\maketitle

\begin{abstract}
  Laporan ini mengkaji sinyal audio dari berkas \texttt{DRUMS.mp3} dalam
  paket proyek \textit{Bad Piggies Theme Remake} yang dirujukkan pada video
  Rick Ruyz. Audio dapat didekode langsung dengan Python sehingga analisis
  tidak memerlukan FL Studio. Sinyal stereo berdurasi sekitar $160{,}77$
  detik dan berfrekuensi sampling $44\,100$ Hz diubah menjadi sinyal mono,
  kemudian dianalisis menggunakan FFT global dan transformasi Fourier
  waktu-singkat (STFT). Notasi dan normalisasi mengikuti bahan kuliah
  Analisis Fourier dan Wavelet. Dua jendela Hann dengan panjang $1024$
  dan $4096$ sampel digunakan untuk membandingkan lokalisasi waktu dan
  frekuensi. Analisis mencakup puncak spektrum, distribusi energi,
  perubahan karakteristik lokal, pemeriksaan Parseval, dan hubungan
  lebar jendela dengan Prinsip Ketaktentuan Heisenberg. Puncak terpilih
  terkuat berada pada $41{,}60$ Hz dan sekitar $91{,}19\%$ energi berada
  pada pita $0$--$250$ Hz. Perubahan RMS serta distribusi daya lokal
  menunjukkan karakteristik nonstasioner secara deskriptif. Cakupan hasil
  adalah berkas audio pendamping yang tersedia, bukan hasil render
  keseluruhan aransemen proyek FLP.
\end{abstract}
\noindent\textbf{Kata kunci:} transformasi Fourier, FFT, STFT, spektrogram,
jendela Hann, audio.

\section{Pendahuluan}
Audio musik tersusun atas osilasi, perubahan amplitudo, serta kejadian
singkat yang dapat memiliki kandungan frekuensi berbeda. Grafik waktu
memperlihatkan perubahan amplitudo, sedangkan representasi Fourier
memperlihatkan kontribusi komponen frekuensi. Akan tetapi, magnitudo
spektrum global tidak langsung menunjukkan kapan suatu komponen
menjadi kuat. Kebutuhan akan informasi lokal tersebut menjadi alasan
penggunaan STFT, sebagaimana dibahas pada bahan kuliah keenam
\cite{kuliah6}.

Objek yang digunakan adalah \texttt{DRUMS.mp3}, salah satu berkas
audio pendamping dalam folder
\texttt{[ FREE FLP ] Bad Piggies - Theme NEW REMAKE FL STUDIO}.
Menurut keterangan penyedia data untuk proyek ini, paket tersebut
dirujukkan pada video Rick Ruyz berjudul
\textit{[ FREE FLP ] Bad Piggies - Theme REMAKE FL STUDIO STOCK PLUGINS}
\cite{youtube}. Judul video dan nama kanal diperiksa melalui metadata
publik YouTube; kesamaan audio lokal dengan soundtrack video tidak
diuji. Penamaan \texttt{DRUMS} dipakai sebagai identitas berkas, bukan
sebagai bukti bahwa seluruh komponennya berasal dari satu instrumen.

Format FLP menyimpan informasi proyek FL Studio dan menggunakan
sampel pendamping \cite{flp}. Dalam pekerjaan ini, berkas FLP tidak
dirender. Berkas MP3 yang sudah tersedia dapat dibaca langsung sebagai
deret sampel sehingga kebutuhan analisis Fourier tetap terpenuhi tanpa
aplikasi tersebut. Berkas \texttt{VOCALS.mp3} dan sampel
\texttt{Classical Acoustic} tidak dijumlahkan dengan data utama karena
posisi, pengaturan instrumen, dan efek aransemen tidak direkonstruksi.

Tujuan proyek adalah menghubungkan representasi waktu, spektrum global,
dan spektrogram melalui argumentasi matematis serta hasil komputasi
yang dapat dijalankan ulang. Secara khusus, dua lebar jendela dibandingkan
untuk memperlihatkan kompromi lokalisasi waktu dan frekuensi pada
sinyal audio yang sama.

\section{Dasar Teori}
\subsection{Notasi dan domain pengamatan}
Notasi dasar mengikuti Kuliah 4--6 pada folder \texttt{Materi}
\cite{kuliah4,kuliah5,kuliah6}. Simbol $f$ digunakan untuk sinyal;
frekuensi dalam Hz dilambangkan dengan $\nu$ agar berbeda dari sinyal.
Tabel~\ref{tab:notasi} merangkum simbol yang digunakan.

\begin{table}[H]
  \centering\small
  \caption{Simbol dan pengertiannya.}\label{tab:notasi}
  \begin{tabular}{p{0.23\linewidth}p{0.69\linewidth}}
    \toprule
    Simbol                       & Pengertian                                                                  \\
    \midrule
    $f(t)$, $\widehat f(\omega)$ & Sinyal waktu dan transformasi Fourier uniter.                               \\
    $t$, $\tau$                  & Waktu pengamatan dan pusat lokalisasi STFT, dalam detik.                    \\
    $f_s$, $\Delta t$            & Frekuensi sampling (Hz) dan interval sampling $1/f_s$ (detik).              \\
    $N$, $T$                     & Banyak sampel dan durasi pengamatan $T=N\Delta t$.                          \\
    $t_k$, $f_k$                 & Waktu $k\Delta t$ dan nilai sinyal pada indeks $k$.                         \\
    $\nu$, $\omega$              & Frekuensi (Hz) dan frekuensi sudut (rad/s), $\omega=2\pi\nu$.               \\
    $F_m$                        & Koefisien DFT pada indeks frekuensi $m$.                                    \\
    $w$, $w_r$                   & Jendela kontinu dan sampel jendela pada indeks lokal $r$.                   \\
    $M$, $H$, $q$                & Panjang jendela, langkah pergeseran dalam sampel, dan indeks bingkai.       \\
    $Z_{q,m}$                    & DFT segmen berjendela pada bingkai $q$.                                     \\
    $V_wf$, $\Spec_wf$           & STFT dan kuadrat magnitudonya (spektrogram).                                \\
    $P_m$, $P_{q,m}$             & Kerapatan daya satu sisi, masing-masing global dan lokal, per Hz.           \\
    $c_m$                        & Bobot satu sisi: $2$ pada pasangan frekuensi interior, $1$ pada DC/Nyquist. \\
    $\sigma_t$, $\sigma_\omega$  & Simpangan baku waktu dan frekuensi sudut berbobot energi jendela.           \\
    \bottomrule
  \end{tabular}
\end{table}

\subsection{Transformasi Fourier dan energi}
Untuk sinyal yang memenuhi syarat integrabilitas, konvensi dalam
Kuliah 4, Bagian 4.2, adalah
\begin{equation}
  \widehat f(\omega)=\frac{1}{\sqrt{2\pi}}
  \int_{\R}f(t)e^{-i\omega t}\,dt,
  \qquad
  f(t)=\frac{1}{\sqrt{2\pi}}\int_{\R}\widehat f(\omega)e^{i\omega t}\,d\omega.
  \label{eq:fourier}
\end{equation}
Untuk $f\in L^2(\R)$, transformasi dapat diperluas dalam pengertian
$L^2$. Teorema Plancherel memberikan
\begin{equation}
  E=\|f\|_2^2=\int_{\R}|f(t)|^2\,dt
  =\int_{\R}|\widehat f(\omega)|^2\,d\omega.
  \label{eq:plancherel}
\end{equation}
Dengan demikian, energi dapat dihitung melalui domain waktu maupun
domain frekuensi \cite{kuliah4}. Pada audio digital, amplitudo yang
digunakan tidak dikalibrasi sebagai tekanan akustik. Oleh sebab itu,
energi dalam laporan ini merupakan energi sinyal digital, bukan energi
fisik dalam joule.

\subsection{Sampling, DFT, dan FFT}
Pengamatan diskret dinyatakan sebagai
\begin{equation}
  t_k=k\Delta t,\quad f_k=f(t_k),\quad k=0,\ldots,N-1,
  \qquad \Delta t=\frac1{f_s},\quad T=N\Delta t.
\end{equation}
Sampel terakhir terletak pada $(N-1)\Delta t$, sedangkan $T$ menyatakan
durasi blok berisi $N$ sampel. Batas Nyquist adalah $f_s/2$; pengamatan
diskret tidak menentukan secara unik komponen analog di atas batas
tersebut tanpa asumsi tambahan \cite{kuliah5}.

DFT dan inversnya mengikuti Kuliah 5, Bagian 5.3--5.4:
\begin{equation}
  F_m=\sum_{k=0}^{N-1}f_k e^{-2\pi i mk/N},\qquad
  f_k=\frac1N\sum_{m=0}^{N-1}F_m e^{2\pi i mk/N}.
  \label{eq:dft}
\end{equation}
FFT merupakan algoritma untuk menghitung DFT tersebut. Normalisasi
maju yang tidak memuat faktor $1/N$ sesuai dengan konvensi bawaan
implementasi FFT \cite{numpy}. Pada sisi frekuensi nonnegatif,
\begin{equation}
  \nu_m=\frac{mf_s}{N},\quad
  \omega_m=2\pi\nu_m,\quad
  \Delta\nu=\frac{f_s}{N}=\frac1T,\quad
  \widehat f(\omega_m)\approx\frac{\Delta t}{\sqrt{2\pi}}F_m.
  \label{eq:konversi}
\end{equation}
Pendekatan terakhir adalah aproksimasi integral untuk sinyal dalam
interval pengamatan yang diperpanjang dengan nol di luar interval.
DFT sendiri memiliki interpretasi periodisasi deret berhingga. Kisi
frekuensi yang rapat tidak menjamin setiap puncak dapat diinterpretasikan
sebagai satu nada yang menetap sepanjang rekaman.

Kesamaan Parseval diskret memberikan pemeriksaan yang tepat secara
aljabar:
\begin{equation}
  E_{\mathrm{waktu}}=\Delta t\sum_{k=0}^{N-1}|f_k|^2
  =\frac{\Delta t}{N}\sum_{m=0}^{N-1}|F_m|^2
  =E_{\mathrm{frekuensi}}.
  \label{eq:parseval}
\end{equation}
Untuk sinyal real, pasangan frekuensi negatif merupakan konjugat
pasangan positifnya. Karena itu, spektrum satu sisi dapat digunakan
dengan bobot $c_m$ untuk mempertahankan energi. Jika $N$ ganjil,
tidak ada bin tepat di Nyquist dan seluruh bin positif berbobot $2$.
Kerapatan daya global didefinisikan sebagai
\begin{equation}
  P_m=\frac{c_m|F_m|^2}{f_sN},\qquad
  \sum_{m\geq0}P_m\Delta\nu=\frac1N\sum_k|f_k|^2.
  \label{eq:psdglobal}
\end{equation}
Jadi integral diskret $P_m$ memberikan rerata kuadrat, dan hasilnya
dikalikan $T$ untuk memperoleh energi. Tinggi puncak kerapatan daya
berbeda makna dengan persentase energi suatu pita frekuensi.

\subsection{STFT, jendela Hann, dan spektrogram}
Definisi STFT dalam Kuliah 6, Bagian 6.3, adalah
\begin{equation}
  V_wf(\tau,\omega)=\frac1{\sqrt{2\pi}}
  \int_{\R} f(t)\overline{w(t-\tau)}e^{-i\omega t}\,dt.
  \label{eq:stft}
\end{equation}
Spektrogramnya adalah
\begin{equation}
  \Spec_wf(\tau,\omega)=|V_wf(\tau,\omega)|^2.
  \label{eq:spec}
\end{equation}
STFT menganalisis $f$ setelah dilokalisasi oleh jendela. Untuk jendela
real, konjugat dalam~\eqref{eq:stft} tidak mengubah nilainya.
Perkalian oleh jendela di domain waktu bersesuaian dengan konvolusi
spektrum. Akibatnya, spektrum lokal mengandung pengaruh sinyal
sekaligus pelebaran dari respons frekuensi jendela \cite{kuliah6}.

Jendela Hann periodik yang digunakan adalah
\begin{equation}
  w_r=\frac12-\frac12\cos\left(\frac{2\pi r}{M}\right),
  \qquad r=0,\ldots,M-1.
  \label{eq:hann}
\end{equation}
Penyebut $M$ sesuai dengan \texttt{hann(M, sym=False)} untuk analisis
spektral \cite{hann}. Definisi ini merupakan pengambilan sampel jendela
Hann kontinu pada interval setengah terbuka; berbeda dari varian
simetris dengan penyebut $M-1$.

Dengan langkah $H$ dan tanpa penambahan nol pada sinyal, digunakan
bingkai yang memenuhi $qH+M\leq N$:
\begin{equation}
  Z_{q,m}=\sum_{r=0}^{M-1} f_{qH+r}w_r e^{-2\pi i mr/M},\quad
  \tau_q=\frac{qH+M/2}{f_s},\quad \nu_m=\frac{mf_s}{M}.
  \label{eq:stftdiskret}
\end{equation}
Hubungannya dengan konvensi kontinu adalah
\begin{align}
  V_wf(\tau_q,\omega_m)     & \approx
  \frac{\Delta t}{\sqrt{2\pi}}
  e^{-i\omega_m qH\Delta t}Z_{q,m},   \\
  \Spec_wf(\tau_q,\omega_m) & \approx
  \frac{\Delta t^2}{2\pi}|Z_{q,m}|^2.
  \label{eq:speckonversi}
\end{align}
Faktor fase akibat awal bingkai tidak memengaruhi kuadrat magnitudo.
Untuk membandingkan dua panjang jendela secara konsisten, grafik
menampilkan kerapatan daya satu sisi
\begin{equation}
  P_{q,m}=\frac{c_m|Z_{q,m}|^2}{f_s\sum_r w_r^2},\qquad
  \Spec_wf\approx P_{q,m}\frac{\sum_r w_r^2}{2\pi f_s c_m}.
  \label{eq:psdstft}
\end{equation}
Grafik tersebut merupakan spektrogram yang diskalakan sebagai
kerapatan daya, bukan nilai mentah $\Spec_wf$. Penjumlahan frekuensinya
memenuhi
\begin{equation}
  \sum_{m\geq0}P_{q,m}\frac{f_s}{M}
  =\frac{\sum_r|f_{qH+r}w_r|^2}{\sum_r w_r^2},
  \label{eq:parsevalframe}
\end{equation}
yang merupakan rerata kuadrat lokal berbobot energi jendela.
Warna dinyatakan sebagai $10\log_{10}(P_{q,m}/P_0)$,
dengan $P_0=1$ satuan amplitudo digital kuadrat per Hz. Acuan dan
batas warna dibuat sama untuk kedua jendela. Nilai ini bukan tingkat
tekanan suara dalam dB SPL.

\subsection{Resolusi dan Prinsip Ketaktentuan Heisenberg}
Durasi jendela $a=M/f_s$, jarak antarpusat bingkai $H/f_s$, dan jarak
bin frekuensi $f_s/M$ adalah tiga besaran yang berbeda. Jendela lebih
panjang menghasilkan jarak bin lebih kecil, tetapi setiap koefisien
melibatkan rentang waktu lebih panjang. Pada Hann periodik,
\begin{equation}
  B_{\mathrm{ENBW}}=f_s\frac{\sum_r w_r^2}{(\sum_r w_r)^2}
  =\frac32\frac{f_s}{M},
  \label{eq:enbw}
\end{equation}
sedangkan lebar lobus utama antara dua nol pertama sekitar $4f_s/M$.
Besaran tersebut membantu membedakan jarak kisi dari kemampuan
memisahkan komponen yang berdekatan.

Heisenberg dalam materi kuliah dinyatakan menggunakan simpangan baku
berbobot energi. Untuk $w\ne0$, tetapkan
\begin{align}
  t_w             & =\frac{\int t|w(t)|^2dt}{\|w\|_2^2},                                 &
  \omega_w        & =\frac{\int\omega|\widehat w(\omega)|^2d\omega}{\|\widehat w\|_2^2},   \\
  \sigma_t^2      & =\frac{\int(t-t_w)^2|w(t)|^2dt}{\|w\|_2^2},                          &
  \sigma_\omega^2 & =\frac{\int(\omega-\omega_w)^2|\widehat w(\omega)|^2d\omega}
                     {\|\widehat w\|_2^2}.
\end{align}
Jika kedua variansi berhingga, berlaku \cite{kuliah6}
\begin{equation}
  \sigma_t\sigma_\omega\geq\frac12,
  \qquad \sigma_t\sigma_\nu\geq\frac1{4\pi}.
  \label{eq:heisenberg}
\end{equation}
Dengan demikian, perkalian $a(f_s/M)=1$ bukan pemeriksaan langsung
ketaksamaan tersebut: $a$ dan $f_s/M$ bukan simpangan baku dalam
definisinya.

\section{Permasalahan}
Domain pengamatan adalah seluruh sinyal mono yang didekode dari
\texttt{DRUMS.mp3}. Pertanyaan yang dijawab adalah sebagai berikut.
\begin{enumerate}
  \item Bagaimana perubahan amplitudo sinyal terhadap waktu, dan bagaimana
        frekuensi serta interval sampling menentukan representasi diskretnya?
  \item Frekuensi apa yang menghasilkan puncak spektrum global, dan
        bagaimana kontribusi energi terbagi ke dalam pita frekuensi?
  \item Bagaimana distribusi frekuensi lokal berubah terhadap waktu;
        apakah rekaman menunjukkan karakteristik nonstasioner dalam
        pengertian sinyal deterministik?
  \item Bagaimana perbedaan spektrogram Hann dengan $M=1024$ dan
        $M=4096$ sampel, serta hubungannya dengan Heisenberg?
  \item Informasi apa yang dapat diperoleh dari spektrum global dan
        STFT, serta apa keterbatasan penafsirannya untuk berkas ini?
\end{enumerate}

\section{Pembahasan}
\subsection{Data dan alur komputasi}
Data dibaca dari berkas MP3 lokal menggunakan Python. Tidak dilakukan
pengunduhan ulang audio YouTube ataupun konversi FLP menjadi audio.
Sumber video dicantumkan sebagai rujukan asal paket sesuai keterangan
pengguna. Identitas berkas, hash SHA-256, versi pustaka, dan parameter
komputasi dicatat bersama hasil agar objek yang dianalisis dapat
diperiksa kembali.

Kedua kanal stereo dirata-ratakan, kemudian rerata globalnya dikurangi:
\begin{equation}
  g_k=\frac{L_k+R_k}{2},\quad
  \overline g=\frac1N\sum_{k=0}^{N-1}g_k,\quad
  f_k=g_k-\overline g.
  \label{eq:mono}
\end{equation}
Sinyal $f_k$ inilah yang digunakan dalam FFT maupun STFT. Pengurangan
rerata menghilangkan offset DC; tidak dilakukan perubahan frekuensi
sampling atau normalisasi amplitudo masing-masing segmen. Perataan
kanal dapat mereduksi komponen yang berlawanan fase, sehingga hasil
berlaku untuk proyeksi mono ini dan tidak mewakili seluruh informasi
spasial stereo.

\begin{table}[H]
  \centering
  \caption{Parameter data setelah dekode dan prapemrosesan.}\label{tab:data}
  \begin{tabular}{lr}
    \toprule
    Besaran                                   & Nilai                \\
    \midrule
    Kanal sumber                              & 2 (stereo)           \\
    Jumlah sampel tiap kanal, $N$             & $7\,090\,031$        \\
    Frekuensi sampling, $f_s$                 & $44\,100$ Hz         \\
    Interval sampling, $\Delta t$             & $22{,}675737$ $\mu$s \\
    Durasi, $T$                               & $160{,}771678$ s     \\
    Batas Nyquist                             & $22\,050$ Hz         \\
    Rerata mono yang dikurangi, $\overline g$ & ${,}000149161$       \\
    RMS global setelah pengurangan rerata     & ${,}077565858$       \\
    Amplitudo absolut maksimum mono           & ${,}892528833$       \\
    \bottomrule
  \end{tabular}
\end{table}
Alur komputasinya adalah dekode audio, pembentukan $f_k$, penggambaran
waktu dan RMS, FFT global, perhitungan STFT untuk dua jendela,
perbandingan fitur lokal, serta ekspor grafik dan tabel. Seluruh hasil
dalam laporan berasal dari eksekusi notebook, dengan Python 3.11.5,
NumPy 1.26.4, SciPy 1.16.2, dan SoundFile 0.13.1. Kedua analisis
memakai frekuensi sampling asli dan data mono yang sama.

\subsection{Representasi dalam domain waktu}
Grafik waktu pada Gambar~\ref{fig:waktu} memperlihatkan sinyal sepanjang
rekaman serta cuplikan yang lebih singkat. Kurva RMS digunakan sebagai
ringkasan kekuatan amplitudo lokal. Untuk suatu blok indeks $I_j$,
\begin{equation}
  \RMS_j=\left(\frac1{|I_j|}\sum_{k\in I_j}|f_k|^2\right)^{1/2}.
\end{equation}
RMS mengabaikan tanda amplitudo dan lebih mudah dibandingkan
antarbagian daripada tinggi satu sampel ekstrem.
\begin{figure}[H]
  \centering\includegraphics[width=\linewidth]{01_waktu_rms.png}
  \caption{Representasi waktu dan RMS audio utama setelah konversi mono.
    Cuplikan waktu memperjelas osilasi yang tidak terlihat pada skala
    seluruh rekaman.}\label{fig:waktu}
\end{figure}
RMS lokal dihitung pada blok $4410$ sampel ($100$ ms), dengan
penyesuaian jumlah sampel pada blok terakhir. Nilai RMS terbesar
sekitar $0,22455$ terdapat pada blok berpusat di $t=84{,}85$ s.
Ringkasan min--maks setiap $10$ ms pada panel pertama mempertahankan
informasi amplitudo ekstrem tanpa menggambar jutaan titik sekaligus.
Pada cuplikan $40$--$43$ s tampak kejadian amplitudo besar yang
berulang, diselingi amplitudo lebih kecil. Pola ini memotivasi
pemeriksaan apakah kejadian tersebut juga mengandung perubahan
distribusi frekuensi, bukan hanya perubahan tingkat amplitudo.

Nilai RMS menunjukkan seberapa kuat amplitudo sinyal pada suatu bagian rekaman. Sebagai gambaran, ketukan drum yang kuat seperti "DUM!" biasanya menghasilkan amplitudo efektif lebih besar daripada ketukan yang pelan.
Pada penelitian ini, RMS lokal terbesar ditemukan di sekitar detik 84,85. Artinya, bagian tersebut mempunyai amplitudo efektif yang relatif tinggi.

\subsection{Analisis Fourier global}
FFT diterapkan pada seluruh $N$ sampel tanpa jendela tambahan, sehingga
pengamatan berhingga secara implisit memakai jendela persegi panjang.
Spektrum yang dihasilkan menggambarkan keseluruhan interval dan dapat
memuat kebocoran spektral dari batas pengamatan serta perubahan
komponen di dalamnya. Gambar~\ref{fig:global} menampilkan spektrum satu
sisi dan pembesaran rentang frekuensi rendah.
\begin{figure}[H]
  \centering\includegraphics[width=\linewidth]{02_spektrum_global.png}
  \caption{Kerapatan daya FFT global ditampilkan sebagai selubung
    maksimum kelompok frekuensi $5$ Hz pada panel penuh dan ${,}5$ Hz
    pada pembesaran. Tabel puncak dan energi tetap dihitung dari bin asli.
    Sumbu menyatakan $\nu$ dalam Hz; $\omega=2\pi\nu$.}\label{fig:global}
\end{figure}
Jarak bin FFT seluruh rekaman adalah
$\Delta\nu=0{,}006220001$ Hz. Tabel~\ref{tab:puncak} merangkum delapan
puncak lokal terpilih dalam rentang $20$--$10\,000$ Hz. Jarak minimum
antarpuncak yang dipilih adalah $25$ Hz agar tabel tidak hanya memuat
bin-bin yang berdekatan dalam satu gugus. Nilai dB relatif dihitung
terhadap kerapatan daya puncak terpilih terkuat.
\begin{table}[H]
  \centering
  \caption{Puncak spektrum terpilih; angka tabel menggunakan titik desimal.}
  \label{tab:puncak}
  \input{hasil_numerik/tabel_puncak.tex}
\end{table}
Puncak terkuat dalam pemilihan tersebut berada pada
$\nu\approx41{,}599367$ Hz, atau $\omega\approx261{,}376530$ rad/s.
Puncak berikutnya berada di sekitar $75{,}40$ dan $122{,}20$ Hz.
Kerapatan daya puncak pertama sekitar $0,0333335$ satuan amplitudo
kuadrat per Hz. Nilai ini dapat melebihi rerata kuadrat sinyal karena
merupakan kerapatan per satuan frekuensi, bukan integral seluruh daya.

Fraksi energi pita dihitung melalui
$\eta_B=\sum_{\nu_m\in B}P_m\Delta\nu/\sum_mP_m\Delta\nu$.
Pita menggunakan batas bawah inklusif dan batas atas eksklusif,
kecuali pita terakhir yang mencakup bin tertinggi.
\begin{table}[H]
  \centering
  \caption{Distribusi energi global menurut pita frekuensi.}
  \label{tab:pita}
  \input{hasil_numerik/tabel_pita.tex}
\end{table}
Sekitar $91{,}19\%$ energi berada di bawah $250$ Hz. Energi di atas
$4000$ Hz sekitar $5{,}20\%$, sehingga komponen frekuensi tinggi tetap
terukur meskipun energi global didominasi frekuensi rendah. Persentase
ini dihitung dari seluruh bin asli, bukan dari kurva selubung pada
gambar. Pembulatan tabel dapat menyebabkan jumlahnya sedikit berbeda
dari $100\%$.

Pemeriksaan~\eqref{eq:parseval} menghasilkan
\[
  E_{\mathrm{waktu}}={,}9672767344661772,\qquad
  E_{\mathrm{frekuensi}}={,}9672767344661792,
\]
dalam satuan amplitudo digital kuadrat dikali detik. Galat relatifnya
sekitar $2{,}07\times10^{-15}$, konsisten dengan pembulatan numerik.
Integral kerapatan daya satu sisi memberikan rerata kuadrat
$0,0060164623$, yang sama dengan kuadrat RMS global. Pemeriksaan ini
mendukung konsistensi normalisasi dan pembobotan frekuensi; tidak
membuktikan bahwa seluruh informasi sumber analog atau proyek FLP
telah direkonstruksi.

Hasil FFT memperlihatkan bahwa salah satu puncak frekuensi terkuat berada pada 41,60 Hz. Frekuensi tersebut termasuk frekuensi rendah yang, jika dibayangkan sebagai bunyi, berhubungan dengan getaran lambat dan berat seperti "buuum", bukan bunyi tinggi seperti "ting".

Selain itu, sekitar 91,19\% energi digital rekaman berada pada rentang 0-250 Hz. Artinya, energi sinyal lebih banyak terkonsentrasi pada frekuensi rendah.

\subsection{Analisis waktu--frekuensi dan perubahan lokal}
STFT dihitung dengan dua jendela Hann, yaitu $M=1024$ dan $M=4096$,
serta $H=M/4$ sehingga tumpang tindihnya $75\%$. Panjang FFT per
bingkai sama dengan $M$. Hanya bingkai lengkap yang digunakan;
sampel ujung yang tidak membentuk bingkai lengkap tidak dipaksa masuk
melalui penambahan nol. Karena itu, spektrum global tetap mencakup
seluruh berkas, sedangkan pusat bingkai STFT berhenti sebelum ujungnya.

\begin{figure}[H]
  \centering\includegraphics[width=\linewidth]{03_spektrogram_penuh.png}
  \caption{Spektrogram seluruh rekaman dalam bentuk kerapatan daya
    satu sisi dengan acuan dB dan batas warna yang sama. Untuk tampilan
    penuh, nilai maksimum dalam kelompok bingkai dirangkum menjadi
    sekitar $2000$ kolom; lebar transien pada grafik ini tidak dipakai
    sebagai ukuran resolusi. Tampilan dibatasi sampai $8000$ Hz,
    sedangkan perhitungan mencakup seluruh sisi nonnegatif.}\label{fig:specfull}
\end{figure}

Untuk mengukur perubahan bentuk spektrum selain perubahan amplitudo,
digunakan sentroid spektral dan fraksi daya pita:
\begin{equation}
  C_q=\frac{\sum_m\nu_m P_{q,m}}{\sum_m P_{q,m}},\qquad
  R_{q,B}=\frac{\sum_{\nu_m\in B}P_{q,m}}{\sum_m P_{q,m}}.
  \label{eq:centroid}
\end{equation}
Keduanya hanya ditafsirkan untuk bingkai dengan daya tidak nol.
Sentroid adalah pusat massa distribusi daya terhadap frekuensi,
bukan perkiraan nada dasar. Fraksi pita memungkinkan perubahan
bentuk spektrum dibedakan dari perubahan tingkat daya keseluruhan.
\begin{figure}[H]
  \centering\includegraphics[width=\linewidth]{05_perubahan_temporal.png}
  \caption{Perubahan sentroid dan distribusi daya pita pada STFT
    $M=4096$, diringkas per detik dengan bobot daya bingkai. Indikator
    dihitung dari daya, bukan dari nilai dB.}
  \label{fig:temporal}
\end{figure}
Spektrogram penuh menunjukkan konsentrasi kuat pada frekuensi rendah
dan kejadian yang menyebar ke frekuensi lebih tinggi. Sebagian
kejadian tampak sebagai guratan vertikal: kenaikan energi berlangsung
pada beberapa frekuensi dalam waktu yang berdekatan. Karena grafik
penuh diringkas dalam kelompok bingkai, pembahasan detail lebar
kejadian menggunakan pembesaran pada bagian berikutnya.

Sebagai ringkasan numerik tambahan, rekaman dipisahkan menjadi segmen
tidak tumpang tindih sepanjang $20$ s, dengan segmen terakhir lebih
pendek. Tabel~\ref{tab:segmen} menggunakan periodogram setiap segmen
tanpa jendela tambahan, bukan rerata langsung piksel spektrogram.
Rumus sentroid dan fraksi pitanya sama seperti~\eqref{eq:centroid},
tetapi $P_{q,m}$ diganti periodogram segmen. Dengan demikian, tabel
dan grafik lokal menjelaskan variasi temporal pada skala yang berbeda.
\begin{table}[H]
  \centering\small
  \caption{Ringkasan temporal per segmen; dua pita ditampilkan untuk
    perbandingan. Segmen terakhir hampir senyap.}\label{tab:segmen}
  \resizebox{\linewidth}{!}{\input{hasil_numerik/tabel_segmen.tex}}
\end{table}
Pada $40-60$ s, RMS sekitar $0,10215$, sentroid $558{,}99$ Hz,
dan fraksi pita $0$--$250$ Hz sekitar $92{,}86\%$. Pada $120$--$140$ s,
RMS turun menjadi ${,}05724$, tetapi sentroid naik menjadi
$910{,}64$ Hz dan fraksi pita rendah turun menjadi $79{,}85\%$.
Fraksi pita di atas $4000$ Hz meningkat dari $4{,}81\%$ menjadi
$9{,}64\%$. Perubahan fraksi ini menunjukkan bahwa perbedaannya
tidak cukup dijelaskan dengan satu faktor pengali amplitudo tetap:
perkalian sinyal oleh konstanta akan mengubah seluruh daya dengan
faktor yang sama dan mempertahankan fraksi pitanya.

Bagian $160$ s sampai akhir memiliki RMS hanya sekitar $0,000162$.
Pada bagian hampir senyap ini, pengurangan rerata global dapat
meninggalkan offset lokal yang sebanding dengan amplitudonya.
Sentroid dan fraksi pita segmen tersebut tidak digunakan sebagai
bukti utama perubahan karakter suara.

Ini berarti bagian kedua mempunyai amplitudo efektif yang lebih kecil, tetapi pusat distribusi frekuensinya bergeser ke arah frekuensi tinggi.
Dalam bahasa sehari-hari, kondisi ini bisa dianalogikan dengan bunyi "DUM!" yang berat dibandingkan bunyi "tak" atau "cess" yang lebih tajam. Bunyi yang lebih pelan tidak selalu memiliki frekuensi lebih rendah.
Jadi, hasil ini menunjukkan bahwa perubahan karakter suara tidak hanya berkaitan dengan keras-pelannya sinyal, tetapi juga dengan bagaimana energi tersebar pada berbagai frekuensi.

\subsection{Pengaruh lebar jendela}
Gambar~\ref{fig:zoom} memperlihatkan interval yang sama untuk kedua
jendela. Pembandingan dilakukan pada koordinat waktu, frekuensi,
serta skala warna yang sama agar perbedaan tampilan berasal dari
parameter jendela dan bukan pengaturan grafik.
\begin{figure}[H]
  \centering\includegraphics[width=\linewidth]{04_spektrogram_zoom.png}
  \caption{Pembesaran spektrogram pada interval $40$--$43$ detik.
    Jendela pendek melokalisasi perubahan waktu pada rentang yang lebih
    sempit; jendela panjang menghasilkan struktur frekuensi yang lebih
    sempit dengan rentang pengamatan waktu lebih lebar.}\label{fig:zoom}
\end{figure}
\begin{table}[H]
  \centering\small
  \caption{Parameter kedua jendela Hann periodik.}\label{tab:jendela}
  \input{hasil_numerik/tabel_jendela.tex}
\end{table}
Lebar lobus utama antara dua nol pertama sekitar $172{,}27$ Hz untuk
$M=1024$ dan $43{,}07$ Hz untuk $M=4096$. Pada cuplikan
Gambar~\ref{fig:zoom}, guratan frekuensi pada panel jendela panjang
lebih tipis secara vertikal, sementara awal dan akhir kejadian
tersebar pada rentang waktu yang lebih lebar. Misalnya, kejadian
sekitar $40{,}2$ s dan $41{,}0$ s terlihat lebih terlokalisasi pada
panel jendela pendek. Ini merupakan pengamatan bentuk tampilan;
tidak diartikan sebagai pengukuran durasi asli masing-masing sumber.

Jumlah bingkai adalah $27\,692$ untuk $M=1024$ dan $6920$ untuk
$M=4096$. Sampel ekor yang tidak tercakup bingkai lengkap masing-masing
$111$ sampel ($2{,}52$ ms) dan $879$ sampel ($19{,}93$ ms).
Validasi daya dilakukan pada tiga bingkai representatif untuk setiap
jendela dengan~\eqref{eq:parsevalframe}; galat relatif terbesar
sekitar $2{,}11\times10^{-16}$.

Memperbesar $M$ empat kali memperkecil jarak bin, ENBW, dan lebar
lobus utama menjadi seperempat, sekaligus memperpanjang dukungan
waktu empat kali. Karena $H=M/4$, jarak bingkai juga membesar empat
kali. Mengurangi $H$ saja dapat memperpadat kisi waktu, tetapi tidak
menghilangkan perataan waktu oleh jendela selebar $M/f_s$.
Demikian pula, penambahan nol pada FFT hanya memperpadat kisi
frekuensi tanpa mempersempit respons frekuensi jendela.
\begin{figure}[H]
  \centering\includegraphics[width=\linewidth]{06_jendela.png}
  \caption{Jendela Hann dan respons frekuensinya. Respons jendela
    menjelaskan mengapa perubahan panjang dukungan waktu disertai
    perubahan lebar spektral.}\label{fig:jendela}
\end{figure}

Jika kita ingin mengetahui kapan bunyi drum seperti "TAK!" muncul, potongan audio yang pendek lebih membantu karena perubahan waktu dapat diamati dengan lebih tajam.
Namun, jika kita ingin membedakan dua komponen frekuensi yang berdekatan, potongan audio yang lebih panjang memberikan informasi frekuensi yang lebih rinci.
Karena itu, digunakan jendela 1024 dan 4096 sampel untuk menunjukkan perbedaan kedua kondisi tersebut.
Fungsi Hann digunakan pada kedua jendela untuk mengurangi spectral leakage, yaitu penyebaran energi frekuensi akibat pemotongan sinyal. Fungsi Hann mengecilkan bobot sampel pada bagian tepi potongan secara bertahap.

Hubungan dengan Heisenberg dapat dihitung secara eksplisit pada
model kontinu Hann yang disampel oleh~\eqref{eq:hann}:
\begin{equation}
  w_a(t)=\begin{cases}
    \frac12\bigl(1+\cos(2\pi t/a)\bigr), & |t|\leq a/2, \\
    0,                                   & |t|>a/2.
  \end{cases}
\end{equation}
Karena simetri, $t_w=\omega_w=0$. Integrasi langsung memberikan
$\|w_a\|_2^2=3a/8$ dan
\begin{equation}
  \sigma_t^2=a^2\left(\frac1{12}-\frac5{8\pi^2}\right).
\end{equation}
Sifat turunan Fourier dan Plancherel memberikan
$\int\omega^2|\widehat w_a(\omega)|^2d\omega=\|w_a'\|_2^2$
dengan $\|w_a'\|_2^2=\pi^2/(2a)$, sehingga
\begin{equation}
  \sigma_\omega=\frac{2\pi}{\sqrt3\,a},\qquad
  \sigma_t\sigma_\omega=
  \sqrt{\frac{\pi^2}{9}-\frac56}\approx0{,}513117>\frac12.
  \label{eq:hannheisenberg}
\end{equation}
Untuk $M=1024$ diperoleh $\sigma_t\approx3{,}284$ ms dan
$\sigma_\omega\approx156{,}228$ rad/s; untuk $M=4096$ nilainya
sekitar $13{,}138$ ms dan $39{,}057$ rad/s. Hasil kali tetap, sedangkan
penyebaran waktu dan frekuensi berubah berlawanan arah. Ini merupakan
perhitungan untuk model jendela kontinu, bukan pengukuran
ketidakpastian setiap komponen audio.

\subsection{Hubungan antarhasil dan keterbatasan}
FFT global merangkum kontribusi frekuensi sepanjang $T$, sedangkan
STFT mengulangi analisis tersebut setelah melokalisasi sinyal pada
setiap $\tau_q$. Spektrogram kemudian menampilkan kuadrat magnitudo
STFT. Puncak yang kuat secara global dapat berasal dari kejadian
pendek yang kuat atau komponen lemah yang bertahan lama; spektrogram
membantu membedakan lokalisasi kedua kemungkinan tersebut.

Transformasi Fourier kompleks lengkap tetap dapat diinversikan;
yang tidak langsung tersedia pada grafik magnitudo global adalah
informasi lokasi waktu. Spektrogram juga tidak menampilkan fase
STFT. Karena itu, grafik daya global dan spektrogram tidak identik
dengan seluruh informasi koefisien kompleksnya.

Penjumlahan semua piksel spektrogram secara langsung bukan energi
rekaman: terdapat pembobotan jendela, tumpang tindih, kisi waktu dan
frekuensi, serta cakupan tepi yang perlu diperhitungkan. Dalam
laporan ini, energi global diperiksa menggunakan~\eqref{eq:parseval},
sedangkan konsistensi daya per bingkai diperiksa melalui
\eqref{eq:parsevalframe}.

Keterbatasan hasil mencakup kompresi MP3, pengurangan informasi stereo
ketika kedua kanal dirata-ratakan, dan ketiadaan kalibrasi amplitudo
akustik. Puncak spektrum bukan bukti otomatis identitas instrumen,
partitur, atau frekuensi fundamental. Selain itu, analisis hanya
berlaku untuk \texttt{DRUMS.mp3}; susunan lengkap seluruh instrumen
di dalam FLP tidak direkonstruksi. Jendela tetap juga tidak dapat
sekaligus memberikan lokalisasi waktu sangat tajam dan pemisahan
frekuensi sangat halus untuk semua kejadian dalam rekaman.

\section{Kesimpulan}
Analisis audio dapat dilakukan tanpa FL Studio karena paket proyek
telah menyediakan berkas MP3 yang dapat didekode menjadi sinyal
diskret. FFT global, STFT, dan spektrogram memberikan sudut pandang
yang saling melengkapi terhadap berkas \texttt{DRUMS.mp3}. Konvensi
notasi dan normalisasi kontinu mengikuti bahan kuliah, sedangkan
faktor konversi DFT serta pembobotan satu sisi dijelaskan secara
eksplisit untuk menjaga konsistensi perhitungan energi.
Puncak terpilih terkuat berada pada $41{,}60$ Hz, dan $91{,}19\%$
energi terdapat pada pita $0$--$250$ Hz. Kesamaan energi waktu dan
frekuensi diperiksa dengan galat relatif sekitar
$2{,}07\times10^{-15}$. Perubahan RMS, sentroid, serta fraksi daya
antarsegmen mendukung karakteristik nonstasioner secara deskriptif,
bukan sekadar perubahan amplitudo yang seragam.

Perbandingan Hann $1024$ dan $4096$ sampel menunjukkan kompromi
antara lokalisasi waktu dan frekuensi. Jendela pendek mencakup sekitar
$23{,}22$ ms dengan jarak bin $43{,}07$ Hz, sedangkan jendela panjang
mencakup sekitar $92{,}88$ ms dengan jarak bin $10{,}77$ Hz.
Perubahan tersebut sesuai dengan Prinsip Ketaktentuan Heisenberg;
untuk model kontinu Hann, hasil kali simpangan bakunya sekitar
$0,513117$. Pengamatan lokal diperlukan untuk menafsirkan perubahan
temporal, sementara spektrum global berguna untuk merangkum distribusi
frekuensi keseluruhan. Hasil ini belum menentukan aransemen lengkap,
identitas instrumen secara pasti, ataupun sifat probabilistik suatu
proses pembangkit audio.

\begin{thebibliography}{9}
  \bibitem{kuliah4}
  \textit{Kuliah 4: Transformasi Fourier}. Bahan kuliah Analisis Fourier
  dan Wavelet, tanpa tahun. Berkas \texttt{Kuliah-04-AFW.pdf}, folder
  \texttt{Materi}, khususnya Bagian 4.2, hlm.~52--53, dan Bagian 4.5,
  hlm.~58--59.
  \bibitem{kuliah5}
  \textit{Kuliah 5: Sampling dan Analisis Fourier Diskret}. Bahan kuliah
  Analisis Fourier dan Wavelet, tanpa tahun. Berkas
  \texttt{KuliahAFW-05.pdf}, folder \texttt{Materi}, khususnya Bagian
  5.1, hlm.~65, serta Bagian 5.3--5.5, hlm.~69--75.
  \bibitem{kuliah6}
  \textit{Kuliah 6: Analisis Waktu-Frekuensi}. Bahan kuliah Analisis
  Fourier dan Wavelet, tanpa tahun. Berkas \texttt{KuliahAFW-06.pdf},
  folder \texttt{Materi}. Bagian 6.3--6.7, hlm.~89--100, dan petunjuk
  Tugas \textit{Project-Based} 1, Bagian 6.10, hlm.~104.
  \bibitem{youtube}
  Rick Ruyz. \textit{[ FREE FLP ] Bad Piggies - Theme REMAKE FL STUDIO
    STOCK PLUGINS}. Video YouTube, tanpa tahun.
  \url{https://www.youtube.com/watch?v=DdIcPpjQXV4}.
  Judul dan kanal diverifikasi pada 22 Oktober 2026. Rujukan asal paket
  audio lokal berdasarkan keterangan pengguna.
  \bibitem{flp}
  Image-Line. \textit{FL Studio Project File (*.flp)}.
  Dokumentasi FL Studio.
  \url{https://www.image-line.com/fl-studio-learning/fl-studio-online-manual/html/fformats_open_flp.htm}.
  Diakses 22 Oktober 2026.
  \bibitem{numpy}
  NumPy Developers. \textit{Discrete Fourier Transform}.
  \url{https://numpy.org/doc/stable/reference/routines.fft.html}.
  Diakses 22 Oktober 2026.
  \bibitem{hann}
  SciPy Developers. \textit{scipy.signal.windows.hann}.
  \url{https://docs.scipy.org/doc/scipy/reference/generated/scipy.signal.windows.hann.html}.
  Diakses 22 Oktober 2026.
\end{thebibliography}

% \appendix
% \section{Reproduksi komputasi}
% Notebook \texttt{analisis\_audio.ipynb} memuat penjelasan,
% kode analisis, dan keluaran komputasi. Berkas
% \texttt{analisis\_audio.py} menyediakan fungsi analisis yang digunakan
% notebook. Grafik disimpan di \texttt{Gambar}, sedangkan ringkasan
% numerik dan tabel CSV disimpan di \texttt{hasil\_numerik}.
% Jalankan seluruh sel notebook dari awal dengan folder paket audio
% tetap berada di sebelah notebook. Petunjuk instalasi dependensi
% dan pembuatan ulang keluaran tersedia di \texttt{README.md}.
% Notebook tidak memerlukan FL Studio dan tidak mengakses jaringan
% ketika membaca audio lokal.

\end{document}
